\documentclass[a4paper,11pt,twoside]{article}
\usepackage[utf8]{inputenc}
\usepackage{xparse}
\usepackage[english,greek]{babel}
\usepackage{alphabeta}
\usepackage{nimbusserif}
\usepackage[T1]{fontenc}
\usepackage[outer=1.50cm, inner=2.00cm, top=2.00cm, bottom=2.00cm]{geometry}
\usepackage{multicol,longtable,multirow,hhline,enumitem,tikz,pgfplots,tkz-euclide,tkz-tab,capt-of,fontawesome5,gensymb,tabularray,fancyhdr,etoolbox,eurosym,xcolor-material,siunitx,eqparbox,microtype}
\usetikzlibrary{arrows.meta}
\usepackage[most]{tcolorbox}
\usetikzlibrary{tikzmark}
\let\myBbbk\Bbbk
\let\Bbbk\relax
\usepackage[curlybraces,amsbb,mtphrb]{mtpro2}
\usepackage[explicit]{titlesec}
\usepackage{soul}
\newcommand{\eng}{\selectlanguage{english}}
\newcommand{\gr}{\selectlanguage{greek}}
\usepackage{mathimatika}
\def\xrwma{red!80!black}
\definecolor{steelblue}{cmyk}{.7,.278,0,.294}
\definecolor{doc}{cmyk}{1,0.455,0,0.569}
\definecolor{orange}{HTML}{ff7300}

\newcommand{\ekthetesdeiktes}{\DeclareMathSizes{10.95}{10.95}{7}{5}
\DeclareMathSizes{6}{6}{3.8}{2.7}
\DeclareMathSizes{8}{8}{5.1}{3.6}
\DeclareMathSizes{9}{9}{5.8}{4.1}
\DeclareMathSizes{10}{10}{6.4}{4.5}
\DeclareMathSizes{12}{12}{7.7}{5.5}
\DeclareMathSizes{14.4}{14.4}{9.2}{6.5}
\DeclareMathSizes{17.28}{17.28}{11}{7.9}
\DeclareMathSizes{20.74}{20.74}{13.3}{9.4}
\DeclareMathSizes{24.88}{24.88}{16}{11.3}

\makeatletter
\newcommand{\subsup}{
\AtBeginDocument{
\check@mathfonts
\fontdimen16\textfont2=2.5pt
\fontdimen17\textfont2=2.5pt
\fontdimen14\textfont2=4.5pt
\fontdimen13\textfont2=4.5pt}
}
\makeatother}
\usepackage{wrapfig}
\newenvironment{WrapText1}[3][r]
{\wrapfigure[#2]{#1}{#3}}
{\endwrapfigure}

\newenvironment{WrapText2}[3][l]
{\wrapfigure[#2]{#1}{#3}}
{\endwrapfigure}
\newcommand{\wrapr}[6]{
\begin{minipage}{\linewidth}\mbox{}\\
\vspace{#1}
\begin{WrapText1}{#2}{#3}
\vspace{#4}#5\end{WrapText1}#6
\end{minipage}}

\newcommand{\wrapl}[6]{
\begin{minipage}{\linewidth}\mbox{}\\
\vspace{#1}
\begin{WrapText2}{#2}{#3}
\vspace{#4}#5\end{WrapText2}#6
\end{minipage}}
\usepackage{etoolbox,hhline,moreenum}
\usepackage{caption} 
\captionsetup[table]{skip=5pt}
\makeatletter
\newif\ifLT@nocaption
\preto\longtable{\LT@nocaptiontrue}
\appto\endlongtable{%
\ifLT@nocaption
\addtocounter{table}{\m@ne}%
\fi}
\preto\LT@caption{%
\noalign{\global\LT@nocaptionfalse}}
\makeatother

\newlist{alist}{enumerate}{1}
\setlist[alist]{label=\let\textdexiakeraia\relax\alph*.}
\UseTblrLibrary{counter}

\pgfmathdeclarefunction{gauss}{2}{%
\pgfmathparse{1/(#2*sqrt(2*pi))*exp(-((x-#1)^2)/(2*#2^2))}%
}
\pgfkeys{/pgfplots/aks_on/.style={axis lines=center,
xlabel style={at={(current axis.right of origin)},xshift=1.5ex,anchor=center},
ylabel style={at={(current axis.above origin)},yshift=1.5ex, anchor=center}}}
\pgfkeys{/pgfplots/grafikh parastash/.style={red!80!black,line width=.4mm,samples=200}}
\pgfkeys{/pgfplots/belh ar/.style={tick label style={font=\scriptsize},axis line style={-latex}}}
\tikzstyle{pl}=[line width=0.3mm]
\tikzstyle{plm}=[line width=0.4mm]

\newcommand{\kerkissans}[1]{{\fontfamily{maksf}\selectfont {#1}}}
\AtBeginDocument{\renewcommand{\textstigma}{\textsigma\texttau}}

\definecolor{titlecolor}{HTML}{cd0f00}

\newbox\TitleUnderlineTestBox
\newcommand*\TitleUnderline[1]
{%
\bgroup
\setbox\TitleUnderlineTestBox\hbox{\colorbox{titleblue}\strut}%
\setul{\dimexpr\dp\TitleUnderlineTestBox-.3ex\relax}{.3ex}%
\ul{#1}%
\egroup
}
\newcommand*\SectionNumberBox[1]
{%
\colorbox{red!80!black}
{%
\makebox[2em][c]
{%
\color{white}%
\strut
\csname the#1\endcsname
}%
}%
\TitleUnderline{\ \ \ }%
}
\titleformat{\section}[hang]
{\Large\fontfamily{maksf}\selectfont}%
{\colorbox{\xrwma}{%
\raisebox{0pt}[13pt][3pt]{\makebox[80pt]{% height, width
\color{white}{\kerkissans{\textbf{\thesection ο Κεφάλαιο}}}}%
}}}%
{0pt}%
{\colorbox{black}{\raisebox{0pt}[13pt][3pt]{\color{white}\ \textbf{#1}\ }}}

\titleformat{\subsection}[hang]
{\large\bfseries\fontfamily{maksf}\selectfont}%
{\colorbox{red!80!black}{%
\raisebox{0pt}[13pt][3pt]{\makebox[30pt]{% height, width
\color{white}{\kerkissans{\textbf{\thesubsection}}}}%
}}}%
{0pt}%
{
\ \textbf{#1}\ }

\makeatletter
\@addtoreset{section}{part}
\makeatother

\titleformat{\part}[display]
{\normalfont\huge\filcenter\bfseries}{}{-30pt}{\Huge \textcolor{red!80!black}{ \kerkissans{ #1}}}
\titlespacing*{\part} 
{0pt}{0pt}{0pt}

\setlist[enumerate]{itemsep=0mm,label=\textcolor{\xrwma}{\textbf{\thesection.\arabic*}}}
\definecolor{bblue}{HTML}{4F81BD}
\definecolor{rred}{HTML}{C0504D}
\definecolor{ggreen}{HTML}{9BBB59}
\definecolor{ppurple}{HTML}{9F4C7C}


\pgfdeclarelayer{background}
\pgfdeclarelayer{foreground}
\pgfsetlayers{background,main,foreground}

\pgfplotsset{every axis/.append style={
x tick label style={/pgf/number format/.cd, 1000 sep={.}}}}

\fancyhf{}
\newcommand{\myleftmark}{\leftmark}
\renewcommand{\myleftmark}{{\large Τυπολόγιο}}
\renewcommand{\headrulewidth}{1pt}
\renewcommand{\sectionmark}[1]{\markboth{\large Κεφάλαιο \thesection\ -\ #1}{} }

\makeatletter% so we can use macros with @ in their names
\ifthenelse{\boolean{@twoside}}{%
\fancyhead[LE]{
\begin{tikzpicture}[overlay, remember picture]%
    \node[anchor=north west, text=red4, font=\large\scshape, minimum size=1in, inner xsep=5mm] at ($(current page.north west)+(10mm,5mm)$) {\kerkissans{\textbf{\thepage\ $\lvert$ \leftmark}}};
\node[anchor=north east, text=black, font=\large\scshape, minimum size=1in, inner xsep=5mm] at ($(current page.north east)+(-15mm,5mm)$) {\kerkissans{\textbf{\myleftmark}}};
\end{tikzpicture}
}%
\fancyhead[RO]{
\begin{tikzpicture}[overlay, remember picture]%
    \node[anchor=north west, text=red4, font=\large\scshape, minimum size=1in, inner xsep=5mm] at ($(current page.north west)+(15mm,5mm)$) {\kerkissans{\textbf{\myleftmark}}};
\node[anchor=north east, text=black, font=\scshape, minimum size=1in, inner xsep=5mm] at ($(current page.north east)+(-10mm,5mm)$) {\kerkissans{\textbf{\leftmark\ $ \lvert $\ \thepage}}};
\end{tikzpicture}
}%
}
\makeatother

%-------- ΠΑΡΑΤΗΡΗΣΕΙΣ -----------------
\newcounter{parathrhsh}[section]
\renewcommand{\theparathrhsh}{\arabic{parathrhsh}}  
\newcommand{\Parathrhsh}[1]{\refstepcounter{parathrhsh}{\textbf{\textcolor{white}{\faLightbulb}\ \ -\ \ \kerkissans{Παρατήρηση\hspace{2mm}\theparathrhsh}}}}{}

%----------- ΑΣΚΗΣΗ ------------------
\newcommand\thema{A}

\newcounter{askhsh}[section]
\renewcommand{\theaskhsh}{\arabic{askhsh}}   
\newcommand{\Askhsh}{\refstepcounter{askhsh}{\textbf{\textcolor{\xrwma}{\faPenSquare\ \  \large \kerkissans{\textbf{\theaskhshο\hspace{2mm}Θέμα}}}}}\hspace{0mm}}{}

\newlist{erwthma}{enumerate}{3}
\setlist[erwthma]{label=\bf{\kerkissans{\large{\textcolor{\xrwma}{\thema.\arabic*}}}},itemsep=0mm,leftmargin=0.8cm}
%------------------------------------
\newenvironment{askhsh}[1]
{\begin{tcolorbox}[title=\Askhsh\kerkissans{\textbf{\large{  #1}}},breakable,
enhanced standard,titlerule=-.2pt,toprule=0pt, rightrule=0pt, bottomrule=0pt,
colback=white,left=2mm,top=1mm,bottom=0mm,
boxrule=0pt,
colframe=white,borderline west={1.5mm}{0pt}{\xrwma},leftrule=2mm,sharp corners,coltitle=\xrwma]
\renewcommand{\thema}{#1}}
{\end{tcolorbox}}

\begin{document}
\NineColors{saturation=high}
\pagestyle{fancy}
\pagenumbering{arabic}

\renewcommand{\myleftmark}{\large Τυπολόγιο}
\begin{center}
\part{Επιπλέον Επαναληπτικά Θέματα Β (Μέρος 2ο)}
\end{center}
\section{}

\begin{askhsh}{B}
Δίνονται οι συναρτήσεις $f(x) = e^x - 1$ και $g(x) = \ln x$. 
\begin{erwthma}
\item Να βρείτε το πεδίο ορισμού και τον τύπο των συναρτήσεων $h = \dfrac{f}{g}$ και $k = f \cdot g$.
\item Να υπολογίσετε το όριο $\lim\limits_{x \to 1} \dfrac{f(x)}{g(x)}$.
\item Να αποδείξετε ότι η εξίσωση $k(x) = 1$ έχει τουλάχιστον μία πραγματική ρίζα στο διάστημα $(1, e)$.
\item Να βρείτε τις ασύμπτωτες της γραφικής παράστασης της συνάρτησης $f$.
\end{erwthma}
\end{askhsh}

\begin{askhsh}{B}
Δίνεται η συνάρτηση $f(x) = x \ln x$ με πεδίο ορισμού $(0, +\infty)$.
\begin{erwthma}
\item Να υπολογίσετε το όριο $\lim\limits_{x \to 0^+} f(x)$ και στη συνέχεια, να βρείτε την κατακόρυφη ασύμπτωτη της $C_f$, αν υπάρχει.
\item Να μελετήσετε την $f$ ως προς τη μονοτονία και την κυρτότητα.
\item Να αποδείξετε ότι υπάρχει τουλάχιστον ένα $\xi \in (1, e)$ τέτοιο ώστε $f'(\xi) = \dfrac{e}{e-1}$.
\item Να υπολογίσετε το εμβαδόν του χωρίου που περικλείεται από την $C_f$, τον άξονα $x'x$ και τις ευθείες $x=1$ και $x=e$.
\end{erwthma}
\end{askhsh}

\begin{askhsh}{B}
Δίνεται η συνάρτηση $f(x) = x(x-2)e^x$ με πεδίο ορισμού το $\mathbb{R}$.
\begin{erwthma}
\item Να αποδείξετε ότι η εξίσωση $x^2 - 2 = 0$ έχει τουλάχιστον μία λύση στο $(0, 2)$.
\item Να υπολογίσετε το όριο $\lim\limits_{x \to -\infty} f(x)$.
\item Να βρείτε, αν υπάρχει, την οριζόντια ασύμπτωτη της γραφικής παράστασης της $f$.
\item Να μελετήσετε την $f$ ως προς τη μονοτονία και να βρείτε τα τοπικά ακρότατα.
\end{erwthma}
\end{askhsh}

\begin{askhsh}{B}
Δίνεται η συνάρτηση $f(x) = \dfrac{e^x - x - 1}{x^2}$ για $x \ne 0$.
\begin{erwthma}
\item Να υπολογίσετε το όριο $\lim\limits_{x \to 0} f(x)$.
\item Να ορίσετε τον αριθμό $\alpha \in \mathbb{R}$ ώστε η συνάρτηση $g(x) = \begin{cases} f(x), & x \ne 0 \\ \alpha, & x = 0 \end{cases}$ να είναι συνεχής στο $\mathbb{R}$.
\item Να αποδείξετε ότι υπάρχει $x_0 \in (0, 1)$ τέτοιο ώστε η συνεχής συνάρτηση $g$ να παίρνει την τιμή $g(x_0) = 0,7$. Δίνεται ότι $e \approx 2,71$.
\item Να αποδείξετε ότι $\lim\limits_{x \to +\infty} \dfrac{\ln x}{f(x)} = 0$.
\end{erwthma}
\end{askhsh}

\begin{askhsh}{B}
Δίνεται η συνάρτηση $f(x) = x^3 - 3x + 1$ με $x \in \mathbb{R}$.
\begin{erwthma}
\item Να αποδείξετε ότι η εξίσωση $f(x)=0$ έχει τουλάχιστον μία ρίζα $x_0 \in (0, 1)$.
\item Να μελετήσετε την $f$ ως προς τη μονοτονία.
\item Να εξηγήσετε γιατί το προηγούμενο $x_0$ είναι μοναδικό στο διάστημα $(0, 1)$.
\item Να αποδείξετε ότι υπάρχει σημείο της γραφικής παράστασης της $f$ στο οποίο η εφαπτομένη είναι παράλληλη στο τμήμα που ενώνει τα σημεία $A(0, f(0))$ και $B(2, f(2))$ και να βρείτε την τετμημένη του $\xi \in (0, 2)$.
\end{erwthma}
\end{askhsh}

\begin{askhsh}{B}
Δίνονται οι συναρτήσεις $f(x) = \ln(x^2+1)$ και $g(x) = x$, ορισμένες στο $\mathbb{R}$.
\begin{erwthma}
\item Να ορίσετε τη συνάρτηση $h(x) = f(x) - g(x)$ και να αποδείξετε ότι το όριο $\lim\limits_{x \to 0} \dfrac{h(x)}{x^2}$ ισούται με $1$.
\item Να αποδείξετε τη σχέση $\ln(x^2+1) \le x$ για κάθε $x \ge 0$. Πότε ισχύει η ισότητα?
\item Να μελετήσετε τη συνάρτηση $h$ ως προς τη μονοτονία.
\item Να βρείτε, αν υπάρχει, την πλάγια ασύμπτωτη της γραφικής παράστασης της $h$ στο $+\infty$.
\end{erwthma}
\end{askhsh}

\begin{askhsh}{B}
Δίνεται η συνάρτηση $f(x) = \dfrac{\ln x}{x-1}$ για $x > 0$ και $x \ne 1$.
\begin{erwthma}
\item Να εξετάσετε αν η συνάρτηση
\[g(x)=\begin{cdcases}
f(x) &, 0<x\neq 1\\ 1 & , x=1
\end{cdcases}\]
είναι συνεχής.
\item Δίνεται η συνάρτηση $h(x) = (x-1)^2f(x)$. Να δείξετε ότι υπάρχει ρίζα της εξίσωσης $h(x) = 1$ στο διάστημα $(1, e)$.
\item Να βρείτε τις ασύμπτωτες της $C_f$.
\end{erwthma}
\end{askhsh}

\begin{askhsh}{B}
Δίνονται οι συναρτήσεις $f(x) = x+1$ και $g(x) = \dfrac{1}{x+1}$.
\begin{erwthma}
\item Να βρείτε το πεδίο ορισμού και τον τύπο της συνάρτησης $h = f + g$.
\item Να εξετάσετε αν η εξίσωση $h(x)=0$ έχει λύση στο διάστημα $[-2, 0]$ και αν υπάρχει $x_0 \in [0, 1]$ τέτοιο ώστε $h(x_0)=2$.
\item Να μελετήσετε την $h$ ως προς τη μονοτονία και να βρείτε τα τοπικά ακρότατα.
\item Να βρείτε τις ασύμπτωτες της γραφικής παράστασης της $h$.
\end{erwthma}
\end{askhsh}

\begin{askhsh}{B}
Δίνεται η συνάρτηση $f(x) = e^x - x - 1$ με $x \in \mathbb{R}$.
\begin{erwthma}
\item Να βρείτε το όριο $\lim\limits_{x \to 0} \dfrac{f(x)}{x^2}$.
\item Να αποδείξετε ότι $f(x) \ge 0$ για κάθε $x \in \mathbb{R}$.
\item Να εξετάσετε τη συνάρτηση ως προς την κυρτότητα.
\item Να δείξετε ότι η εξίσωση $f(x) = 2$ έχει ακριβώς μία ρίζα στο διάστημα $(1, 2)$. (Δίνεται $e \approx 2,7, e^2 \approx 7,4$).
\end{erwthma}
\end{askhsh}

\begin{askhsh}{B}
Δίνεται η συνάρτηση $f(x) = x e^x$ με πεδίο ορισμού το $\mathbb{R}$.
\begin{erwthma}
\item Να υπολογίσετε το $\lim\limits_{x \to -\infty} f(x)$.
\item Να συμπεράνετε ποια είναι η οριζόντια ασύμπτωτη στο $-\infty$.
\item Να αποδείξετε ότι η γραφική παράσταση της $f$ τέμνει την ευθεία $y = -1$ σε ένα τουλάχιστον σημείο με την τετμημένη του στο διάστημα $(-1, 0)$.
\item Να προσδιορίσετε το ολικό ελάχιστο της συνάρτησης.
\end{erwthma}
\end{askhsh}

\begin{askhsh}{B}
Δίνεται η συνάρτηση $f(x) = \dfrac{\sqrt{x+1}-1}{x}$ για $x \ge -1, x \ne 0$.
\begin{erwthma}
\item Να υπολογίσετε το $\lim\limits_{x \to 0} f(x)$.
\item Έστω $g$ η συνεχής επέκταση της $f$ στο $[-1, +\infty)$. Να βρείτε τον τύπο της.
\item Να αποδείξετε ότι η συνεχής συνάρτηση $g$ λαμβάνει την τιμή $\dfrac{2}{5}$ για τουλάχιστον ένα $\xi \in (0, 3)$.
\item Να αποδείξετε ότι η $g$ αντιστρέφεται στο πεδίο ορισμού της.
\end{erwthma}
\end{askhsh}

\begin{askhsh}{B}
Δίνεται η πολυωνυμική συνάρτηση $f(x) = x^4 - 2x^2 + 1$.
\begin{erwthma}
\item Να αποδείξετε ότι υπάρχει τουλάχιστον ένα $\xi \in (-1, 1)$ στο οποίο η εφαπτομένη της γραφικής παράστασης της $f$ είναι οριζόντια.
\item Να βρείτε τις τιμές $\xi \in (-1, 1)$ για τις οποίες μηδενίζεται η παράγωγος $f'(\xi) = 0$.
\item Να μελετήσετε τη συνάρτηση ως προς τη μονοτονία και να εντοπίσετε τα τοπικά της ακρότατα.
\item Να βρείτε τα σημεία καμπής της $C_f$.
\end{erwthma}
\end{askhsh}

\begin{askhsh}{B}
Δίνεται η συνάρτηση $f(x) = e^x (x-1) + 1$, με $x \in \mathbb{R}$.
\begin{erwthma}
\item Να υπολογίσετε το $\lim\limits_{x \to -\infty} f(x)$.
\item Να μελετήσετε την $f$ ως προς τη μονοτονία και την κυρτότητα.
\item Να αποδείξετε ότι υπάρχει μοναδικό $x_0 \in (0, 1)$ τέτοιο ώστε $f(x_0) = \dfrac{1}{2}$.
\item Να βρείτε την εξίσωση της εφαπτομένης της $C_f$ στο σημείο της $M(0, f(0))$.
\end{erwthma}
\end{askhsh}

\begin{askhsh}{B}
Δίνεται η συνάρτηση $f(x) = \dfrac{x - \hm x}{x + \hm x}$.
\begin{erwthma}
\item Να βρείτε το $\lim\limits_{x \to 0} f(x)$.
\item Να υπολογίσετε το $\lim\limits_{x \to +\infty} f(x)$.
\item Να αποδείξετε ότι η συνάρτηση $h(x) = x - \hm x$ είναι γνησίως αύξουσα.
\item Να δικαιολογήσετε γιατί η εξίσωση $x = \hm x$ έχει μοναδική ρίζα σε όλο το $\mathbb{R}$.
\end{erwthma}
\end{askhsh}

\begin{askhsh}{B}
Δίνεται η συνάρτηση $f(x) = \ln x$ με πεδίο ορισμού $(0, +\infty)$.
\begin{erwthma}
\item Να αποδείξετε ότι υπάρχει $\xi \in (x, x+1)$ τέτοιο ώστε $\ln(x+1) - \ln x = \dfrac{1}{\xi}$.
\item Να αποδείξετε την ανισότητα $\dfrac{1}{x+1} < \ln\left(1 + \dfrac{1}{x}\right) < \dfrac{1}{x}$ για κάθε $x > 0$.
\item Με βάση την παραπάνω ανισότητα, αποδείξτε ότι $\lim\limits_{x \to +\infty} \left[x \ln\left(1 + \dfrac{1}{x}\right)\right] = 1$.
\item Να επαληθεύσετε το όριο του ερωτήματος Β3 με δεύτερο τρόπο.
\end{erwthma}
\end{askhsh}

\begin{askhsh}{B}
Δίνεται η συνάρτηση $f(x) = e^{-x^2}$.
\begin{erwthma}
\item Να αποδείξετε ότι υπάρχει $\xi \in (-1, 1)$ στο οποίο η εφαπτομένη της γραφικής παράστασης της $f$ είναι οριζόντια και να το βρείτε.
\item Να μελετήσετε την $f$ ως προς τη μονοτονία και να βρείτε το μέγιστο της συνάρτησης.
\item Να βρείτε τα σημεία καμπής της γραφικής παράστασης της $f$.
\item Να υπολογίσετε το όριο $\lim\limits_{x \to +\infty} f(x)$ και να γράψετε την αντίστοιχη ασύμπτωτη.
\end{erwthma}
\end{askhsh}

\begin{askhsh}{B}
Δίνεται η συνάρτηση $f(x) = \dfrac{e^x}{x}$ για $x > 0$.
\begin{erwthma}
\item Να αποδείξετε ότι υπάρχει τουλάχιστον ένα $x_0 \in (1, 2)$ τέτοιο ώστε $3f(x_0) = e^2+e$.
\item Να ελέγξετε αν η εξίσωση $3f(x) = e^2+e$ έχει πάνω από μία λύση στο $(1, 2)$, μελετώντας τη μονοτονία της $f$.
\item Να βρείτε το $\lim\limits_{x \to +\infty} f(x)$.
\item Να βρείτε την εξίσωση της εφαπτομένης της $C_f$ η οποία διέρχεται από την αρχή των αξόνων $O(0,0)$.
\end{erwthma}
\end{askhsh}

\begin{askhsh}{B}
Δίνεται η συνάρτηση $f(x) = x^3 + x - 1$.
\begin{erwthma}
\item Να αποδείξετε ότι υπάρχει τουλάχιστον ένα $\xi \in (0, 1)$ τέτοιο ώστε η εφαπτομένη της $C_f$ στο $M(\xi, f(\xi))$ να έχει κλίση $1$. Βρείτε το ακριβές $\xi$.
\item Να αποδείξετε ότι η εξίσωση $f(x) = 0$ έχει τουλάχιστον μία ρίζα στο $(0, 1)$.
\item Να αποδείξετε ότι αυτή η ρίζα είναι μοναδική στο $\mathbb{R}$.
\item Υπολογίστε το εμβαδόν του χωρίου που περικλείεται από τη $C_f$, τον άξονα $x'x$ και τις ευθείες $x=1$, $x=2$.
\end{erwthma}
\end{askhsh}

\begin{askhsh}{B}
Δίνεται η συνάρτηση $f(x) = x^2 \ln x$ με $x > 0$.
\begin{erwthma}
\item Να υπολογίσετε το $\lim\limits_{x \to 0^+} f(x)$.
\item Να βρείτε την τιμή της συνεχούς επέκτασης της $f$ στο $0$.
\item Να αποδείξετε ότι υπάρχει $\xi \in (1, e)$ τέτοιο ώστε η εφαπτομένη της $C_f$ στο $\xi$ να είναι παράλληλη στο τμήμα με άκρα $A(1, f(1))$ και $B(e, f(e))$.
\item Να μελετήσετε την κυρτότητα της συνάρτησης.
\end{erwthma}
\end{askhsh}

\begin{askhsh}{B}
Δίνονται οι συναρτήσεις $f(x) = e^x + 2x$ και $g(x) = x^2 + 1$.
\begin{erwthma}
\item Να ορίσετε τον τύπο της συνάρτησης σύνθεσης $f \circ g$.
\item Να αποδείξετε ότι υπάρχει ένα τουλάχιστον $x_0 \in (0, 1)$ τέτοιο ώστε $f(x_0) = 3$.
\item Να αποδείξετε ότι η εξίσωση $f'(x) = 0$ δεν έχει λύση στο $\mathbb{R}$. Ποιο είναι το συμπέρασμα για τη μονοτονία της $f$;
\item Να δείξετε ότι υπάρχει $\xi \in (0, x_0)$ ώστε $f'(\xi) = \dfrac{2}{x_0}$.
\end{erwthma}
\end{askhsh}

\end{document}
